\section{An exact choice of weights for Gaussian Hecke characters}
Let $K=\mathbb Q(i)$. Every nonzero ideal has a generator $a$, and
$\eta((a))=(a/|a|)^4$ is independent of the choice of generator.
Write $L_k(s)=L_K(s,\eta^k)$ and $L_0=\zeta_K$.
For $k\ne0$ these are entire Hecke functions of nonzero infinity type.
Conjugation of ideals gives $L_{-k}=L_k$.
Their zeros enter, for example, the explicit formula for counts of Gaussian prime angles.
See \href{https://arxiv.org/abs/1705.07498}{Rudnick and Waxman, \emph{Angles of Gaussian primes}}.

For a real even trigonometric polynomial on the unit circle write
$P(z)=h_0+\sum_{k=1}^N h_k(z^k+z^{-k})$.
In the half-plane $\Re s>1$, termwise logarithmic differentiation gives
\[
 D_P(s)=\sum_{\mathfrak p,m\ge1}
 \frac{\log N\mathfrak p\,P(\eta(\mathfrak p)^m)}{(N\mathfrak p)^{ms}}
 =-h_0\frac{\zeta'_K}{\zeta_K}(s)
 -2\sum_{k=1}^N h_k\frac{L'_k}{L_k}(s).
\]
Here $\mathfrak p$ ranges over prime ideals, including the prime above 2. A nonnegative $P$
gives nonnegative local weights. Nonnegative $h_k$ ensure that orders of zeros
add with the same sign. Its mean $h_0$ is the residue of $D_P$ at 1.
Thus one can ask for the smallest main term compatible with prescribed
sensitivity to several $L_k$ at once.

\subsection{An optimal polynomial for a block of modes}
Fix $0<\alpha<1$ and integers $1\le r\le N$. Define rising factorials
$(a)_0=1$, $(a)_j=a(a+1)\cdots(a+j-1)$ and set
\[
 c_j=\frac{(\alpha)_j}{j!},\qquad b_j=\frac{(2\alpha)_j}{j!},\qquad
 d=\min\{r,\lfloor N/2\rfloor+1\}.
\]
\noindent\textbf{Theorem.}
Among nonnegative $P$ of degree at most $N$ satisfying $h_{N-j}\ge b_j$ for
$0\le j<r$, the exact minimum of $h_0$ is
\[
 M_{N,r}(\alpha)=2\sum_{j=0}^{d-1}c_j^2
 -\boldsymbol{1}_{\{N=2d-2\}}c_{d-1}^2.
\]
It is attained with every $h_k\ge0$. An extremizer is
$P(z)=|\sum_{i=0}^N v_i z^i|^2$, where $v_i=c_{\min(i,N-i)}$ if
$\min(i,N-i)<d$ and $v_i=0$ otherwise.

The parameter $\alpha=1/2$ gives equal required gains $b_j=1$. With $r=N$
this is precisely Belov's theorem after multiplying his cosine polynomial by two.
Belov also proved uniqueness for that case. The weighted block and its use in
the Hecke criterion are developed here. No priority claim is made for the general formula.
See \href{https://doi.org/10.1070/IM1994v043n03ABEH001582}
{Belov, \emph{On an extremal problem on the minimum of a trigonometric polynomial}}.

\medskip\noindent\emph{Proof.}
The coefficients $c_i$ decrease and $(1-z)^{-\alpha}=\sum c_i z^i$,
so $\sum_{i=0}^j c_ic_{j-i}=b_j$. By symmetry,
$h_{N-j}=\sum_{i=0}^j v_i v_{j-i}$.
For $j<d$ this equals $b_j$. If $d\le j<r$, the support covers every index
and $v_i\ge c_i$ throughout. The remaining constraints follow. Parseval gives the stated
mean and every autocorrelation $h_k$ is nonnegative.

For the lower bound use the Fej\'er--Riesz factorization
$P(z)=|\sum_{i=0}^N u_i z^i|^2$, padding with zeros when needed.
Put $A_i=|u_i|$ and $B_i=|u_{N-i}|$ for $0\le i<d$.
Set $w_i=1$, except that $w_{d-1}=1/2$ when $N=2d-2$.
Then $\sum w_i(A_i^2+B_i^2)\le h_0$, since the center index counts once.
Write $e_h=[z^h](1-z)^\alpha$, set $y_i=w_ic_i$, and define
\[
 \lambda_j=\sum_{h=0}^{d-1-j}e_h y_{j+h},\qquad
 H_{ij}=\begin{cases}\lambda_{i+j}& i+j<d,\\0&i+j\ge d.\end{cases}
\]
Here $e_0=1$, $e_h<0$ for $h\ge1$, and $\sum_{h=0}^m e_h>0$,
as follows by taking the coefficient at $z^m$ of $(1-z)^{\alpha-1}$.
Since $y_i$ is positive and nonincreasing, $\lambda_j>0$. The convolution
of $c$ with $e$ is $(1,0,\ldots)$, giving
$(Hc)_i=\sum_{j<d-i}c_j\lambda_{i+j}=w_ic_i$.

For any $i,j$, the arithmetic--geometric mean inequality gives
$2 A_i B_j\le A_i^2 c_j/c_i+B_j^2 c_i/c_j$. Multiply by $H_{ij}$
and sum, using symmetry and $Hc=(w_i c_i)_i$. It follows that
\[
 2\sum_{j<d}\lambda_j b_j
 \le 2\sum_{j<d}\lambda_j h_{N-j}
 \le 2\sum_{i,j<d} H_{ij}A_i B_j
 \le \sum_{i<d}w_i(A_i^2+B_i^2)\le h_0.
\]
The first sum on the left, without the factor two, is
$\sum_{i,j<d}H_{ij}c_i c_j=\sum_{i<d}w_ic_i^2$.
This proves sharpness.\hfill$\square$

\medskip
For example, for $N=r=4$ and $\alpha=1/3$ the factor is
$1+z/3+2z^2/9+z^3/3+z^4$, and the optimal mean is $184/81$.
It meets the first three upper-frequency constraints with equality.
The lower-frequency constraints beyond $\lfloor N/2\rfloor+1$ never increase the optimum.
Stirling's formula gives $c_j\sim j^{\alpha-1}/\Gamma(\alpha)$. For the full
block as $N\to\infty$ this yields
\[
 M_{N,N}(\alpha)=
 \begin{cases}
 2\sum_{j\ge0}c_j^2+o(1), &0<\alpha<1/2,\\
 (2/\pi)\log N+O(1), &\alpha=1/2,\\
 \dfrac{2(N/2)^{2\alpha-1}}{(2\alpha-1)\Gamma(\alpha)^2}(1+o(1)),
 &1/2<\alpha<1 .
 \end{cases}
\]
The bounded-cost regime prescribes decreasing gains. Each of the $N$ gains
is strictly positive at every fixed $N$.

\subsection{The arithmetic estimate required by the extremizer}
Set $g_P=\lim_{s\to1}(D_P(s)-h_0/(s-1))-h_0$, and define
\[
 U_P(x)=\sum_{(N\mathfrak p)^m\le x}\!
 \log N\mathfrak p\,P(\eta(\mathfrak p)^m)
 \left(\frac1{(N\mathfrak p)^m}-\frac1x\right)-h_0\log x-g_P.
\]
\noindent\textbf{Criterion.}
Fix the extremal polynomial. An eventual bound $U_P(x)\le C_\epsilon
x^{-1/2+\epsilon}$ for each $0<\epsilon<1/2$ implies RH for
each $L_k$ with $h_k>0$. In particular, $r=N$ detects every
$L_k$ for $1\le k\le N$.

\medskip\noindent\emph{Proof.}
For $\Re s>1$, termwise integration gives
\[
 \int_1^\infty U_P(x)x^{-s}\,dx=
 \frac{D_P(s)}{s(s-1)}-\frac{h_0}{(s-1)^2}-\frac{g_P}{s-1}.
\]
The expression is regular at 1 by construction. At a zero $\rho$
in the open strip, its residue is
$(h_0\operatorname{ord}_\rho\zeta_K+2\sum h_k\operatorname{ord}_\rho L_k)/
[\rho(1-\rho)]$. This is positive for a real zero detected by $P$, and
nonzero at any detected nonreal zero. The Hecke functions have no zeros on
$\Re s=1$.
Set $\beta=1/2+\epsilon$. Above a suitable cutoff the integrand
$C_\epsilon x^{\beta-1}-U_P(x)$ is nonnegative. Its Mellin transform
continues meromorphically into $\Re s>\beta$. A real pole there would
have negative residue. If its abscissa of convergence exceeded $\beta$,
Landau's theorem would give a singularity at that real abscissa, where
a negative residue contradicts nonnegativity when approached from the right.
Thus its integral is holomorphic on $\Re s>\beta$, excluding every detected zero there.
Let $\epsilon$ decrease to zero and apply the Hecke functional equations.\hfill$\square$

The optimum determines the smallest pole term allowed by this class of detector.
The arithmetic estimate still concerns the angular values on individual Gaussian
prime powers. The source's integral theta-lift comparison for a quadratic
character and a congruence class supplies an estimate for different local coefficients.
For the proposed transfer the unresolved task is to realize an angular weight
in that comparison, already for $N=r=1$.
